\documentclass[12pt,a4paper]{amsart}

\usepackage{amsmath,amssymb,amsthm}
\usepackage{mathtools}
\usepackage{hyperref}
\usepackage{geometry}
\geometry{margin=2.5cm}
\usepackage{booktabs}

% -------------------------------------------------------
% Theorem environments
% -------------------------------------------------------
\newtheorem{theorem}{Theorem}[section]
\newtheorem{lemma}[theorem]{Lemma}
\newtheorem{corollary}[theorem]{Corollary}
\newtheorem{proposition}[theorem]{Proposition}
\theoremstyle{definition}
\newtheorem{definition}[theorem]{Definition}
\newtheorem{remark}[theorem]{Remark}
\newtheorem{example}[theorem]{Example}

% -------------------------------------------------------
% Custom commands
% -------------------------------------------------------
\newcommand{\N}{\mathbb{N}}
\newcommand{\Z}{\mathbb{Z}}
\newcommand{\Q}{\mathbb{Q}}
\newcommand{\R}{\mathbb{R}}
\newcommand{\C}{\mathbb{C}}
\newcommand{\eq}[1]{e\!\left(#1\right)}
\DeclareMathOperator{\ord}{ord}
\DeclareMathOperator{\li}{li}
\DeclareMathOperator{\sgn}{sgn}

% -------------------------------------------------------
\begin{document}
% -------------------------------------------------------

\title{The Hardy--Littlewood Circle Method:\\
       Exponential Sums, Major and Minor Arcs}

\author{Michael Fuhrmann}
\address{Specialist Mathematics Project, Build 80}
\email{specialist-maths@localhost}
\date{2026}

\begin{abstract}
The Hardy--Littlewood \emph{circle method} is one of the most powerful
techniques in analytic number theory.  Introduced by Hardy and Ramanujan
(1918) and developed by Hardy and Littlewood (1920--1928) for additive
problems, it reduces the question of representing an integer $n$ as a sum
of a prescribed number of terms to the evaluation of a contour integral
over the unit circle $[0,1)$.  The key idea is to split this interval into
\emph{major arcs} (neighbourhoods of rational numbers with small
denominator, where exponential sums are large and tractable) and
\emph{minor arcs} (where exponential sums are small and can be bounded).
We develop the method from first principles: exponential sums (Weyl sums
and Gauss sums), the Farey dissection, the structure of major arcs and
their contribution via the \emph{singular series} and \emph{singular
integral}, and the Weyl differencing technique for minor arc bounds.
We illustrate the method on Waring's problem and set up the framework
for Vinogradov's three-prime theorem.
\end{abstract}

\maketitle
\tableofcontents

% ================================================================
\section{Introduction}
% ================================================================

A central problem in additive number theory is: \emph{In how many ways
can a positive integer $n$ be written as
\[
  n \;=\; x_1^k + x_2^k + \cdots + x_s^k,
\]}
where $x_i$ are non-negative integers (Waring's problem), or as
\[
  n \;=\; p_1 + p_2 + \cdots + p_r
\]
where $p_i$ are primes (Goldbach-type problems)?

The circle method transforms such questions into integrals.  Let
$r_s(n)$ denote the number of representations of $n$ as a sum of $s$
$k$-th powers.  The generating function (exponential sum)
\[
  f(\alpha) \;=\; \sum_{m=0}^{M} \eq{m^k \alpha},
  \qquad \eq{x} := e^{2\pi i x},
\]
satisfies the key identity
\begin{equation}\label{eq:key}
  r_s(n) \;=\; \int_0^1 f(\alpha)^s\, \eq{-n\alpha}\, d\alpha.
\end{equation}
This follows immediately from the orthogonality of the characters
$\alpha \mapsto \eq{m\alpha}$: the integral picks out exactly the
coefficient of $\eq{n\alpha}$ in $f(\alpha)^s$.

The circle method is the art of evaluating~\eqref{eq:key} by
decomposing $[0,1)$ into regions where $f(\alpha)$ is either large
(and can be computed explicitly) or small (and can be bounded).

\textbf{Notation.}  Throughout, $\eq{x} = e^{2\pi ix}$, $\log = \log_e$,
$p$ always denotes a prime, $\Lambda$ denotes the von Mangoldt function,
and all implicit constants are absolute unless noted.

% ================================================================
\section{Exponential Sums}
\label{sec:expsums}
% ================================================================

\subsection{Weyl sums}

\begin{definition}[Weyl sum]
  For a polynomial $f(x) = \alpha_k x^k + \cdots + \alpha_1 x \in \R[x]$
  and $N \geq 1$, the \emph{Weyl sum} is
  \[
    W(f; N) \;=\; \sum_{n=1}^{N} \eq{f(n)}.
  \]
  In the simplest case $f(x) = \alpha x^k$, we write $W(\alpha; k; N)$.
\end{definition}

The trivial bound is $|W(f; N)| \leq N$.  The key insight is that when
$\alpha$ is ``far from rationals with small denominator'', cancellation
reduces this significantly.

\begin{lemma}[Weyl, 1916: geometric sum]\label{lem:geom}
  For $\alpha \in \R$ and $N \geq 1$:
  \[
    \left|\sum_{n=1}^{N} \eq{n\alpha}\right|
    \;\leq\; \min\!\left(N,\; \frac{1}{2\|\alpha\|}\right),
  \]
  where $\|x\| = \min_{m \in \Z} |x - m|$ denotes the distance to the
  nearest integer.
\end{lemma}

\begin{proof}
  The sum is a geometric series with ratio $\eq{\alpha}$.  If
  $\eq{\alpha} = 1$, i.e.\ $\alpha \in \Z$, then the sum equals $N$.
  Otherwise:
  \[
    \left|\sum_{n=1}^N \eq{n\alpha}\right|
    = \left|\frac{\eq{\alpha} - \eq{(N+1)\alpha}}{1 - \eq{\alpha}}\right|
    \leq \frac{2}{|1 - \eq{\alpha}|}
    = \frac{1}{|\sin(\pi\alpha)|}
    \leq \frac{1}{2\|\alpha\|},
  \]
  using $|\sin(\pi x)| \geq 2\|x\|$.  Taking the minimum with $N$ gives
  the result.
\end{proof}

\subsection{Gauss sums}

\begin{definition}[Complete Gauss sum]
  For integers $a$, $q$ with $q \geq 1$, the \emph{Gauss sum} is
  \[
    G(a, q) \;=\; \sum_{m=1}^{q} \eq{\frac{a m^k}{q}}.
  \]
\end{definition}

\begin{lemma}[Gauss sum bound]\label{lem:gauss}
  For $\gcd(a, q) = 1$:
  \[
    |G(a, q)| \;\leq\; k \cdot q^{1/2 + \varepsilon}
  \]
  for any $\varepsilon > 0$.  More precisely, for $k = 2$:
  $|G(a,q)| \leq \sqrt{q}$ when $q$ is odd and $q \equiv 1 \pmod 4$;
  in general $|G(a,q)| \leq \sqrt{2q}$.
\end{lemma}

These bounds on Gauss sums are central to controlling major arc
integrals, as we shall see in Section~\ref{sec:major}.

\subsection{Exponential sums over primes}

For prime-counting problems, the relevant exponential sum involves the
von Mangoldt function $\Lambda(n) = \log p$ if $n = p^k$ and $0$ otherwise.

\begin{definition}[Prime exponential sum]
  \[
    S(\alpha) \;=\; S(\alpha; N)
    \;=\; \sum_{n \leq N} \Lambda(n)\, \eq{n\alpha}.
  \]
\end{definition}

This sum has magnitude $\ll N$ on the major arcs (via Dirichlet
characters and the prime number theorem for arithmetic progressions)
and magnitude $\ll N^{1-\delta}$ on the minor arcs (via Vinogradov's
estimates or the large sieve).

% ================================================================
\section{The Farey Dissection}
\label{sec:farey}
% ================================================================

The Farey dissection provides a canonical partition of $[0,1)$ into
arcs centred at rational fractions $a/q$ with small denominator.

\begin{definition}[Farey sequence]
  The \emph{Farey sequence} $\mathcal{F}_Q$ of order $Q$ is the
  set of all fractions $a/q$ with $0 \leq a \leq q \leq Q$ and
  $\gcd(a, q) = 1$ (plus $0/1$ and $1/1$), arranged in increasing order.
\end{definition}

\begin{definition}[Farey arc]
  The \emph{Farey arc} associated to $a/q \in \mathcal{F}_Q$ is the
  interval
  \[
    \mathfrak{M}(a, q)
    \;=\; \left[\frac{a}{q} - \frac{1}{qQ'}, \frac{a}{q} + \frac{1}{qQ''}\right],
  \]
  where $Q'$ and $Q''$ are determined by the neighbours of $a/q$ in
  $\mathcal{F}_Q$.  Crucially, consecutive Farey fractions $a/q < a'/q'$
  satisfy $|a'q - aq'| = 1$ (a fundamental property of the Farey sequence).
\end{definition}

\begin{lemma}[Farey partition]\label{lem:farey}
  The Farey arcs of order $Q$ are pairwise disjoint and cover $[0,1)$:
  \[
    [0, 1) \;=\; \bigsqcup_{q \leq Q} \bigsqcup_{\substack{a=0\\ \gcd(a,q)=1}}^{q-1}
                 \mathfrak{M}(a, q).
  \]
  Each arc $\mathfrak{M}(a, q)$ has length $\sim 1/(qQ)$.
\end{lemma}

\begin{proof}
  This is a standard result in the theory of Farey sequences; see
  Hardy--Wright~\cite{HardyWright}, Chapter~3.  The disjointness follows
  from the mediant property: between consecutive Farey fractions $a/q$
  and $a'/q'$, the midpoint is $(a+a')/(q+q')$, and each point is
  claimed by exactly one arc.
\end{proof}

% ================================================================
\section{Major and Minor Arcs}
\label{sec:arcs}
% ================================================================

In practice, one uses a slightly simplified version of the Farey
dissection, with a fixed parameter $Q$ controlling the ``quality'' of
the rational approximation.

\begin{definition}[Major and minor arcs]\label{def:arcs}
  Let $N \geq 1$ and $Q = Q(N)$ (to be chosen later, typically
  $Q = N^\delta$ for small $\delta > 0$, or $Q = N / (\log N)^B$).
  Define:
  \begin{align*}
    \text{Major arcs:} \quad
    \mathfrak{M} &\;=\; \bigcup_{q \leq Q}\, \bigcup_{\substack{a=0\\ \gcd(a,q)=1}}^{q-1}
                         \mathfrak{M}(a, q), \\[4pt]
    \mathfrak{M}(a, q) &\;=\; \left\{\alpha \in [0,1) :
      \left|\alpha - \frac{a}{q}\right| \leq \frac{Q}{N}\right\}, \\[4pt]
    \text{Minor arcs:} \quad
    \mathfrak{m} &\;=\; [0,1) \setminus \mathfrak{M}.
  \end{align*}
\end{definition}

\begin{remark}
  The major arcs $\mathfrak{M}(a,q)$ are pairwise disjoint provided
  $Q^2 \leq N$, which we henceforth assume.  Each arc has length
  $2Q/N$ (roughly), and there are $\sum_{q \leq Q} \phi(q) \sim Q^2/2$
  such arcs.  The total measure of $\mathfrak{M}$ is $\ll Q^3/N$.
\end{remark}

The philosophy of the circle method is:
\begin{enumerate}
  \item On \textbf{major arcs}, $\alpha$ is close to some $a/q$ with
    $q \leq Q$.  Then $f(\alpha)$ can be approximated by an explicit
    sum (involving Gauss sums) times a simple integral.  The
    contribution is the \emph{main term} of $r_s(n)$.
  \item On \textbf{minor arcs}, $\alpha$ is ``far from small-denominator
    rationals''.  Then $|f(\alpha)|$ is small (by Weyl differencing or
    the large sieve), and the contribution is an \emph{error term}.
\end{enumerate}

The threshold $Q$ must be chosen to balance the two contributions:
large enough so that the minor arcs carry only small error, small enough
so that the major arc computation remains tractable.

% ================================================================
\section{Major Arc Analysis}
\label{sec:major}
% ================================================================

We work out the major arc contribution for Waring's problem:
$f(\alpha) = \sum_{m \leq M} \eq{m^k \alpha}$ where $M = N^{1/k}$.

\begin{lemma}[Major arc approximation]\label{lem:major_approx}
  For $\alpha = a/q + \beta$ with $\gcd(a,q) = 1$, $q \leq Q$, and
  $|\beta| \leq Q/N$:
  \[
    f(\alpha) \;=\; \frac{G(a, q)}{q}\, v(\beta) \;+\; O\!\left(q^{1+\varepsilon}\right),
  \]
  where $v(\beta) = \int_0^M \eq{\beta t^k}\, dt$ is the \emph{Weyl integral}
  (a smooth approximation), and $G(a,q)$ is the Gauss sum of
  Definition~2.3.
\end{lemma}

\begin{proof}
  We split the sum $\sum_{m \leq M} \eq{m^k(a/q + \beta)}$ according to
  residues $m \equiv r \pmod q$.  For each residue $r$:
  \[
    \sum_{\substack{m \leq M\\ m \equiv r (q)}} \eq{m^k(a/q + \beta)}
    \;=\; \eq{r^k a/q} \sum_{j : qj+r \leq M} \eq{(qj+r)^k \beta}.
  \]
  Summing over $r = 1, \ldots, q$ and approximating the inner sum by an
  integral (Euler--Maclaurin with error $O(1)$ per term, so $O(q)$ total):
  \[
    f(\alpha) \;=\; \frac{1}{q}\, G(a, q)\, v(\beta) \;+\; O(q).
  \]
  The improved error $O(q^{1+\varepsilon})$ comes from a more careful
  analysis (see Vaughan~\cite{Vaughan1981}).
\end{proof}

\subsection{The singular series}

\begin{definition}[Singular series]\label{def:singular_series}
  For integers $n, k, s$, the \emph{singular series} is
  \[
    \mathfrak{S}(n) \;=\; \sum_{q=1}^{\infty} \sum_{\substack{a=1\\ \gcd(a,q)=1}}^{q}
                          \left(\frac{G(a,q)}{q}\right)^s \eq{-\frac{an}{q}}.
  \]
  By Lemma~\ref{lem:gauss}, this converges absolutely for $s > 2k+1$.
\end{definition}

\begin{remark}
  The singular series $\mathfrak{S}(n)$ has an Euler product decomposition:
  \[
    \mathfrak{S}(n) \;=\; \prod_p \sigma_p(n),
  \]
  where $\sigma_p(n)$ is the $p$-adic density of solutions to
  $x_1^k + \cdots + x_s^k \equiv n \pmod{p^j}$ for large $j$.
  In particular, $\mathfrak{S}(n) > 0$ if and only if the equation
  $x_1^k + \cdots + x_s^k = n$ has $p$-adic solutions for all $p$.
\end{remark}

\subsection{The singular integral}

\begin{definition}[Singular integral]
  \[
    \mathfrak{J}(n) \;=\; \int_{-\infty}^{\infty} v(\beta)^s\, \eq{-n\beta}\, d\beta,
  \]
  where $v(\beta) = \int_0^1 \eq{\beta u^k}\, du$ (after normalising $M=1$,
  the general case follows by scaling).
\end{definition}

\begin{lemma}[Singular integral estimate]\label{lem:sing_int}
  For $s > k$ and $n > 0$:
  \[
    \mathfrak{J}(n) \;\sim\; C_{k,s} \cdot n^{s/k - 1},
  \]
  where $C_{k,s} = \Gamma(1 + 1/k)^s / \Gamma(s/k)$ is an explicit constant
  depending only on $k$ and $s$.
\end{lemma}

\begin{proof}
  By the substitution $\beta \mapsto \beta/n$ and $u \mapsto (n/s)^{1/k} u$,
  the integral $\mathfrak{J}(n)$ scales as $n^{s/k - 1}$.  The constant
  is evaluated by recognising the integrand as related to the Fourier
  transform of $|u|^{1/k}$; see Vaughan~\cite[Chapter~4]{Vaughan1981}.
\end{proof}

\subsection{Major arc total contribution}

\begin{theorem}[Major arc contribution]\label{thm:major}
  With the notation above, for $s > 2k+1$:
  \[
    \int_{\mathfrak{M}} f(\alpha)^s \eq{-n\alpha}\, d\alpha
    \;=\; \mathfrak{S}(n)\, \mathfrak{J}(n)\, N^{s/k - 1}
          \;+\; O\!\left(N^{s/k - 1 - \delta}\right)
  \]
  for some $\delta = \delta(k, s) > 0$.
\end{theorem}

\begin{proof}[Proof sketch]
  On each major arc $\mathfrak{M}(a, q)$, substitute
  $\alpha = a/q + \beta$ and use Lemma~\ref{lem:major_approx}:
  \[
    \int_{\mathfrak{M}(a,q)} f(\alpha)^s \eq{-n\alpha}\, d\alpha
    \;\approx\; \left(\frac{G(a,q)}{q}\right)^s \eq{-an/q}
                \int_{-Q/N}^{Q/N} v(\beta)^s \eq{-n\beta}\, d\beta.
  \]
  Summing over all $(a, q)$ with $q \leq Q$ and extending the $\beta$-integral
  to $\R$ (with negligible error for $Q$ large), we recover
  $\mathfrak{S}(n) \cdot \mathfrak{J}(n) \cdot N^{s/k-1}$.  Error
  terms arise from: (i) the approximation in Lemma~\ref{lem:major_approx},
  contributing $O(N^{s/k - 1 - 1/k})$; (ii) the tail of the singular
  series, contributing $O(Q^{-1+\varepsilon})$.  Choosing $Q$ appropriately
  balances these.
\end{proof}

% ================================================================
\section{Minor Arc Bounds: Weyl's Inequality}
\label{sec:minor}
% ================================================================

On the minor arcs, we need to show $|f(\alpha)|$ is substantially less
than $N$.  The fundamental tool is:

\begin{theorem}[Weyl's inequality]\label{thm:weyl}
  Let $k \geq 2$ and suppose $|\alpha - a/q| \leq 1/q^2$ for some
  integers $a, q$ with $\gcd(a,q) = 1$ and $1 \leq q \leq N^k$.  Then
  \[
    \left|\sum_{n=1}^{N} \eq{n^k \alpha}\right|
    \;\ll\; N^{1+\varepsilon}
             \left(\frac{1}{q} + \frac{1}{N} + \frac{q}{N^k}\right)^{2^{1-k}}.
  \]
\end{theorem}

\begin{proof}[Proof sketch (Weyl differencing)]
  We sketch the argument for $k = 2$.  Squaring the sum:
  \[
    \left|\sum_n \eq{n^2 \alpha}\right|^2
    = \sum_m \sum_n \eq{(n^2 - m^2)\alpha}
    = \sum_h \sum_n \eq{h(2n + h)\alpha},
  \]
  where $h = n - m$.  Fixing $h$ and summing over $n$, the inner sum is
  a geometric series $\sum_n \eq{2hn\alpha}$, bounded by
  $\min(N, \|2h\alpha\|^{-1})$.  Summing over $h \leq N$ and using
  $\|2h\alpha\| \geq 2h\|\alpha\|$ for generic $\alpha$ yields the
  improvement over the trivial bound.  For general $k$, one iterates
  this differencing $k-1$ times; see Hardy--Wright~\cite[Chapter~21]{HardyWright}
  or Vaughan~\cite[Chapter~2]{Vaughan1981}.
\end{proof}

\begin{corollary}[Minor arc bound]\label{cor:minor}
  On the minor arcs, i.e.\ when $\alpha \in \mathfrak{m}$, every $\alpha$
  satisfies $|q\alpha - a| \geq Q/N$ for all $q \leq Q$.  By Dirichlet's
  theorem, there exist $a, q$ with $q \leq N^k/Q$ and
  $|q\alpha - a| \leq Q/N$.  Applying Weyl's inequality:
  \[
    |f(\alpha)| \;\ll\; N^{1 - \sigma(k) + \varepsilon},
  \]
  where $\sigma(k) = 2^{1-k}$ for suitable $Q$.
\end{corollary}

\begin{remark}[Efficiency for large $k$]
  For $k \geq 2$, Weyl's inequality gives a saving of $N^{-\sigma}$ with
  $\sigma = 2^{1-k}$, which becomes very small as $k \to \infty$.
  Vinogradov's mean value theorem (refined by Wooley's efficient
  congruencing) gives exponentially better bounds, reducing $G(k)$ from
  Weyl's bound of $2^k$ to $k(\log k + O(1))$.
\end{remark}

% ================================================================
\section{The Full Circle Method: Putting It Together}
\label{sec:full}
% ================================================================

We can now state the main theorem of the circle method for Waring's problem.

\begin{theorem}[Hardy--Littlewood--Vinogradov]\label{thm:main}
  For $k \geq 2$ and $s \geq s_0(k)$ (a threshold depending on $k$),
  and for all sufficiently large $n$:
  \[
    r_s(n) \;=\; \mathfrak{S}(n)\, \mathfrak{J}(n)\, n^{s/k - 1}
                 \;+\; O\!\left(n^{s/k - 1 - \delta}\right),
  \]
  where $\mathfrak{S}(n) > 0$ for all $n$ (the local solubility condition
  is automatically satisfied for $s \geq C(k)$).  In particular,
  $r_s(n) > 0$ for all sufficiently large $n$.
\end{theorem}

\begin{proof}[Proof outline]
  Split the integral~\eqref{eq:key} as
  \[
    r_s(n) \;=\; \int_{\mathfrak{M}} f^s \eq{-n\cdot} + \int_{\mathfrak{m}} f^s \eq{-n\cdot}.
  \]
  \textbf{Major arc:} By Theorem~\ref{thm:major}, the major arc integral
  equals $\mathfrak{S}(n)\, \mathfrak{J}(n)\, n^{s/k-1} + O(n^{s/k-1-\delta})$.

  \textbf{Minor arc:} On $\mathfrak{m}$, use Corollary~\ref{cor:minor}
  to bound $|f(\alpha)| \leq N^{1-\sigma}$ and apply two of the $s$ factors
  trivially ($\int_0^1 |f|^2 d\alpha \leq N$, by Parseval).  This gives
  \[
    \left|\int_{\mathfrak{m}} f^s \eq{-n\cdot}\right|
    \;\leq\; N^{(s-2)(1-\sigma)} \cdot \int_0^1 |f|^2
    \;\ll\; N^{s-2-(s-2)\sigma + 1}
    \;=\; N^{s-1-(s-2)\sigma}.
  \]
  For $s > 2/\sigma + 1$, this is $o(N^{s/k - 1})$, so the minor arc is
  absorbed into the error term.
\end{proof}

% ================================================================
\section{Application to Goldbach-Type Problems}
\label{sec:goldbach}
% ================================================================

For prime-counting problems, the exponential sum is
$S(\alpha) = \sum_{n \leq N} \Lambda(n) \eq{n\alpha}$.
The representation count for $r$ primes summing to $n$ is:
\[
  R_r(n) \;=\; \sum_{\substack{p_1 + \cdots + p_r = n\\ p_i \leq N}}
               \log p_1 \cdots \log p_r
  \;=\; \int_0^1 S(\alpha)^r \eq{-n\alpha}\, d\alpha.
\]

\begin{lemma}[Major arc for primes]\label{lem:major_prime}
  For $\alpha = a/q + \beta$ with $\gcd(a, q) = 1$ and $q$ small:
  \[
    S(\alpha) \;\approx\; \frac{\mu(q)}{\phi(q)}\, v_N(\beta),
  \]
  where $v_N(\beta) = \int_2^N \eq{t\beta}\, dt / \log t \approx N\cdot \eq{N\beta}$
  and $\mu$ is the M\"obius function.
\end{lemma}

This approximation holds when $\chi_0$ is the principal character and
the prime number theorem for arithmetic progressions $\pmod{q}$ applies.
The $\mu(q)/\phi(q)$ factor encodes the contribution from all characters
mod $q$ (non-principal characters contribute negligibly on major arcs by
orthogonality and non-vanishing of $L(1,\chi)$).

The minor arc bound for prime sums uses the \emph{large sieve inequality}
(see Paper~15) and Vinogradov's lemma:
\[
  \max_{\alpha \in \mathfrak{m}} |S(\alpha)| \;\ll\; N(\log N)^{-B}
\]
for any $B > 0$, provided $Q = (\log N)^A$ for suitable $A = A(B)$.
This is the content of the Bombieri--Vinogradov theorem and underlies
the proof of Vinogradov's three-prime theorem (Paper~18).

% ================================================================
\section{Historical Notes and Extensions}
\label{sec:history}
% ================================================================

\begin{remark}[History]
  The circle method was first used by Hardy and Ramanujan (1918) to derive
  the asymptotic formula for the partition function $p(n)$.  Hardy and
  Littlewood (1920--1923) applied it to Waring's problem (Papers I--IV),
  establishing the asymptotic formula for $r_s(n)$ and obtaining the
  original bound $G(k) \leq (k-2)2^{k-1} + 5$.  Vinogradov (1937)
  dramatically improved the minor arc bounds via his mean value estimates.
  Wooley (1992--2016) developed \emph{efficient congruencing} and proved
  the Vinogradov mean value conjecture for all $k$; independently,
  Bourgain, Demeter, and Guth (2015) proved the same result via
  $\ell^2$-decoupling (the BDG theorem).  Both proofs are complete and
  Wooley's 2016 efficient-congruencing proof covers all $k$, not only $k=3$.
\end{remark}

\begin{remark}[Limitations]
  The circle method has fundamental limitations:
  \begin{itemize}
    \item It requires $s \geq s_0(k)$ ``many'' summands; for binary Goldbach
      ($r = 2$ primes) it gives only conditional results.
    \item The \emph{parity problem} in sieve theory is closely related to why
      $s = 2$ is out of reach: the circle method, used naively, cannot
      distinguish $p + p$ from $p + p \cdot p'$ (products of two primes).
    \item Nevertheless, for $r \geq 3$ primes, the method succeeds
      unconditionally (Vinogradov's theorem, Paper~18).
  \end{itemize}
\end{remark}

% ================================================================
\section{Conclusion}
% ================================================================

The Hardy--Littlewood circle method rests on a single elegant idea: the
integral representation~\eqref{eq:key} of the representation count.
The method's power comes from the systematic exploitation of the
structure of exponential sums near rational numbers:
\begin{enumerate}
  \item Near a rational $a/q$ with small $q$ (\textbf{major arcs}),
    the exponential sum is approximated by a Gauss sum times an integral.
    Summing over all such rationals gives the \emph{singular series}
    $\mathfrak{S}(n)$ and \emph{singular integral} $\mathfrak{J}(n)$.
  \item Away from all small-denominator rationals (\textbf{minor arcs}),
    the exponential sum is small by Weyl's inequality or the large sieve.
    The minor arc integral is an error term.
\end{enumerate}

This framework successfully resolves Waring's problem and ternary Goldbach,
and underpins much of modern analytic number theory.  The papers that follow
(Papers~18--20) apply and extend this method.

% ================================================================
\begin{thebibliography}{9}

\bibitem{HardyRamanujan1918}
G.~H.~Hardy and S.~Ramanujan.
\newblock Asymptotic formulae in combinatory analysis.
\newblock \emph{Proc.\ London Math.\ Soc.\ (2)}, 17:75--115, 1918.

\bibitem{HardyLittlewood1923}
G.~H.~Hardy and J.~E.~Littlewood.
\newblock Some problems of ``Partitio Numerorum''.  III: On the expression of
  a number as a sum of primes.
\newblock \emph{Acta Math.}, 44:1--70, 1923.

\bibitem{HardyWright}
G.~H.~Hardy and E.~M.~Wright.
\newblock \emph{An Introduction to the Theory of Numbers}, 6th~ed.
\newblock Oxford University Press, 2008.

\bibitem{Vaughan1981}
R.~C.~Vaughan.
\newblock \emph{The Hardy--Littlewood Method}, 2nd~ed.
\newblock Cambridge University Press, 1997.

\bibitem{Vinogradov1954}
I.~M.~Vinogradov.
\newblock \emph{The Method of Trigonometrical Sums in the Theory of Numbers}.
\newblock Interscience Publishers, London, 1954.
  (Translated and revised by K.~F.~Roth and A.~Davenport.)

\bibitem{Wooley2016}
T.~D.~Wooley.
\newblock The cubic case of the main conjecture in Vinogradov's mean value theorem.
\newblock \emph{Adv.\ Math.}, 294:532--561, 2016.

\bibitem{IwaniecKowalski}
H.~Iwaniec and E.~Kowalski.
\newblock \emph{Analytic Number Theory}.
\newblock American Mathematical Society, Providence, RI, 2004.

\end{thebibliography}

\end{document}
